\clearpage
\setcounter{page}{1}
\maketitlesupplementary

% ==========================================================================
% Appendix / Supplementary material.
% Compiled together with the main paper; delete these pages from the main
% submission PDF and/or submit separately per the CVPR instructions.
% ==========================================================================

\noindent This supplementary provides the grid construction and its exact
self-test (\cref{sec:supp-grid}), full architecture and training details for
reproducibility (\cref{sec:supp-arch,sec:supp-train}), the decoding procedure
(\cref{sec:supp-decode}), additional results and visualizations
(\cref{sec:supp-results}), and the planned $\SEthree$ extension
(\cref{sec:supp-se3}).

\section{Grid construction and self-test}
\label{sec:supp-grid}

\paragraph{Hopf--HEALPix cells.}
A cell at level $l$ is indexed by a flat integer $c = \textit{pix}\cdot n_r(l) +
\textit{roll}$, where $\textit{pix}\in\{0,\dots,12\,\textit{nside}^2-1\}$ is a
HEALPix pixel in the \textsc{nested} scheme and $\textit{roll}\in\{0,\dots,
n_r(l)-1\}$ with $n_r(l)=6\,\textit{nside}$ and $\textit{nside}=2^{l}$. The
cell-centre rotation is the Hopf lift of the base point
$(\theta,\phi)=\textsc{pix2ang}(\textit{nside},\textit{pix})$ and fibre angle
$\psi=(\textit{roll}+\tfrac12)\cdot 2\pi/n_r(l)$, giving the ZYZ-Euler rotation
$R=R_z(\phi)R_y(\theta)R_z(\psi)$. The inverse (rotation $\rightarrow$ cell) reads
the base point from $R\hat{z}$ and recovers $\psi$ after removing it; both maps are
exact inverses on cell centres.

\paragraph{Strict child structure.}
In the \textsc{nested} scheme a pixel's four children are $\textit{pix}\cdot4+\{0,
1,2,3\}$ and a fibre bin's two children are $\textit{roll}\cdot2+\{0,1\}$, so a
cell's eight children are indexed by $\textit{child}=\textit{s2child}\cdot2+
\textit{rollchild}$. Hence the path $(c_0,c_1,\dots,c_L)$ with $c_0\in\{0,\dots,
71\}$ and $c_l\in\{0,\dots,7\}$ uniquely identifies a finest cell, and the parent
and the eight children of any node are computed in closed form.

\paragraph{Self-test.}
The whole likelihood path is verified independent of any training: (i)~the
cell$\rightarrow$centre$\rightarrow$cell round trip is exact at every level;
(ii)~decomposing a flat cell into $(c_0,\text{children})$ and recomposing returns
the same cell; and (iii)~a uniform predictor ($\pi_0$ and all $\pi_n$ uniform)
yields average log-likelihood exactly $-\ln\Vso \approx -2.2899$ nats at every
level $0$--$L$, matching \cref{eq:density}. \TODO{optionally report the measured
maximum deviation of the uniform self-test.}

\section{Architecture details}
\label{sec:supp-arch}

\paragraph{Backbone.}
ViT-B/16 (\texttt{vit\_base\_patch16\_224.\allowbreak augreg2\_in21k\_ft\_in1k})~\cite{dosovitskiy2021vit,wightman2019timm},
$224\times224$ input, ImageNet mean/std normalization. We take the
\texttt{forward\_features} patch tokens and drop the prefix (CLS/register) tokens,
giving $P{=}196$ tokens of dimension $D{=}768$.

\paragraph{Cross-attention blocks.}
Both the level-0 scorer and the shared child scorer use the same block: queries
are layer-normed and cross-attend (8 heads) to layer-normed patch tokens
(key/value), followed by a residual MLP with hidden width $4D$ and GELU; two such
layers; a final LayerNorm and linear-to-$1$ projection. The projection weight and
bias are zero-initialized, so each scorer outputs the uniform log-categorical at
initialization.

\paragraph{Queries.}
The level-0 scorer uses $72$ learned query embeddings. The child scorer builds its
$8$ queries from, for each child $j$: the child's offset from the parent centre as
a rotation vector $\in\RR^3$ normalized by $(\Vso/N(l))^{1/3}$, and the flattened
parent rotation $\in\RR^9$; each is Fourier-encoded with $n_\text{freq}{=}3$
(frequencies $\pi\cdot2^k$, sines and cosines, plus the raw value), concatenated,
linearly projected to $\RR^{D}$ through a 2-layer GELU MLP, and summed with a
learned per-level embedding.

\paragraph{Parameter count.}
The ViT-B/16 backbone accounts for $85.8$M parameters and the cross-attention
tree head for $29.1$M, for $114.9$M in total. The head count is independent of the
grid depth $L$, since a single scorer is shared across all nodes and levels.

\section{Training details}
\label{sec:supp-train}

\begin{table}[h]
\centering
\caption{Training hyperparameters.}
\label{tab:supp-hparams}
\small
\begin{tabular}{l l}
\toprule
optimizer & AdamW~\cite{loshchilov2019adamw}, weight decay $0.05$ \\
learning rate (head / backbone) & $10^{-4}$ / $10^{-5}$ \\
schedule & one-cycle~\cite{smith2019onecycle}, $5\%$ warmup \\
batch size & $64$ \\
training steps & $60\mathrm{k}$ \\
gradient clip & $1.0$ \\
precision & fp32 (gradient checkpointing) \\
augmentation & in-plane roll (\cref{sec:aug}) \\
level $L$ (main/abl.) & $6$/$5$ \\
beam width (eval) & $8$ \\
\bottomrule
\end{tabular}
\end{table}

Training minimizes the single-rotation NLL of \cref{eq:loss}; no symmetry
annotations and no equivalent sets are used at training time. The data pipeline
precomputes, per sample, the ground-truth path and the per-child geometry (offsets
and parent rotations) so the network never calls the grid library during training.
Dumping the dataset at level $6$ and slicing paths to the desired $L$ exploits the
prefix property of the tree, so all depth ablations reuse one dump. All runs use a
single RTX~6000~Ada (48\,GB) shared with other resident jobs; the headline $L{=}6$
model completes its $60$k steps (plus the final test eval and per-level decompose)
in about $12$ hours end to end, i.e.\ $\approx\!1.4$ steps/s at batch size $64$.

\section{Decoding}
\label{sec:supp-decode}

Beam-MAP proceeds top-down: keep the top-$b$ level-0 cells by $\log\pi_0$; at each
level expand every kept cell to its $8$ children, score them with the shared child
scorer, form cumulative log-probabilities, and keep the top-$b$ cells overall;
after level $L$, return the cell-centre rotation of the highest-scoring leaf. Total
node scorings $O(bL)$. \TODO{report sensitivity of pose accuracy to $b$
(cross-ref \cref{tab:abl-misc}).}

\section{Additional results}
\label{sec:supp-results}

\paragraph{Per-shape learning curves.}
\TODO{figure: training NLL and held-out average log-likelihood vs.\ steps, per
shape; expected to show the discrete high-symmetry shapes lagging.}

\paragraph{Qualitative distributions.}
\TODO{figure: predicted $\SOthree$ distributions (Hopf/Mollweide visualization)
for one example per shape, with the ground-truth symmetry set overlaid;
contrast cone/cylinder (continuous) with icosahedron (60 discrete modes).}

\paragraph{Per-level breakdown.}
\cref{tab:supp-perlevel} reports the average log-likelihood decomposed by the
contribution of each tree level (how much log-probability each level adds along the
ground-truth path, on $500$ test images/shape from the headline $L{=}6$ model), to
localize where mass is lost. The path log-likelihood is the row sum plus the
density constant $\ln(N/\Vso)=14.46$. Two patterns stand out. The coarsest level
$\ell_0$ carries by far the largest penalty for the high-symmetry shapes
($-4.00$ for icosa, $-3.23$ for cube) because the top-level categorical must divide
mass among the many occupied cells of a large symmetry set, whereas the continuous
shapes pay least there ($-1.80$/$-2.50$ for cone/cyl). At the \emph{finest} levels,
icosa again pays the most ($-1.21$, $-1.57$ at $\ell_5,\ell_6$ versus $-0.73$,
$-1.06$ for tet and $-0.54$, $-0.90$ for cube), i.e.\ its per-mode peaks stay broad
down to the grid floor---the same effect that drives its weaker point accuracy
(\cref{sec:pose}).

\begin{table}[h]
\centering
\caption{Per-level contribution to the ground-truth path log-probability
(SYMSOL~I, headline $L{=}6$ model, $500$ images/shape). $\ell_0$ is the coarsest
categorical; $\ell_1\!-\!\ell_6$ are the per-child categoricals. Row sum plus
$\ln(N/\Vso)=14.46$ equals the path log-likelihood.}
\label{tab:supp-perlevel}
\small
\setlength{\tabcolsep}{4pt}
\begin{tabular}{l ccccccc}
\toprule
Shape & $\ell_0$ & $\ell_1$ & $\ell_2$ & $\ell_3$ & $\ell_4$ & $\ell_5$ & $\ell_6$ \\
\midrule
tet.  & $-2.54$ & $-0.12$ & $-0.11$ & $-0.28$ & $-0.49$ & $-0.73$ & $-1.06$ \\
cube  & $-3.23$ & $-0.11$ & $-0.07$ & $-0.19$ & $-0.37$ & $-0.54$ & $-0.90$ \\
ico.  & $-4.00$ & $-0.44$ & $-0.15$ & $-0.44$ & $-0.90$ & $-1.21$ & $-1.57$ \\
cone  & $-1.80$ & $-0.71$ & $-0.72$ & $-0.74$ & $-0.80$ & $-0.93$ & $-1.20$ \\
cyl.  & $-2.50$ & $-0.71$ & $-0.72$ & $-0.73$ & $-0.78$ & $-0.91$ & $-1.03$ \\
\bottomrule
\end{tabular}
\end{table}

\paragraph{SYMSOL~II data note.}
The raw \texttt{symmetric\_solids} \texttt{tfds} release (v1.0.0) already contains
all eight shapes---the three marked shapes sphereX, cylO and tetX in addition to
the five SYMSOL~I shapes---with both the single \texttt{rotation} label and the
full \texttt{rotations\_equivalent} set (train $8\times45\mathrm{k}$, test
$8\times5\mathrm{k}$). Only our preprocessed path-dump was restricted to the five
SYMSOL~I shapes. The three marked shapes have now been re-dumped and prepared at
$L{=}6$ through the same path-precompute pipeline
(\texttt{data/symsol\_II\_\{tetX,cylO,sphereX\}}, $45$k/$5$k train/test each);
producing \cref{tab:symsol2} needs the three per-shape training runs (one model per
marked shape, the prior-work protocol). Unlike SYMSOL~I, each marked-shape test
image in this release carries a single ground-truth rotation---we verified
$K{=}1$ for all $5{,}000$ test images of every marked shape (SYMSOL~I has
$K{=}12$--$720$)---so the SYMSOL~II metric is log-likelihood at that single pose
and no equivalent set enters evaluation.
\TODO{Document the exact split and seed used for the per-shape runs.}

\subsection{Additional ablations}
\label{sec:supp-ablations}

Beyond the depth, scorer-sharing/beam, and training-length ablations in the main
paper (\cref{sec:ablations}), we set up three further ablations on SYMSOL~I,
reported here as the runs complete.

\paragraph{Backbone.}
To isolate the distribution formulation from the encoder, we match the backbone to
the baselines (\cref{tab:abl-backbone}); head and protocol are held fixed.

\begin{table}[h]
\centering
\caption{\textbf{Ablation: backbone.} Overall SYMSOL~I average log-likelihood with
different image encoders (head and protocol fixed).}
\label{tab:abl-backbone}
\small
\begin{tabular}{l c}
\toprule
Backbone & avg.\ log-lik. \\
\midrule
ResNet-50 & \pend \\
ViT-B/16 (in21k$\rightarrow$in1k) & \pend \\
\bottomrule
\end{tabular}
\end{table}

\paragraph{Child-scorer conditioning.}
The child scorer is conditioned on the normalized child offset, the Fourier-encoded
parent rotation, and a level embedding. \cref{tab:abl-cond} removes each in turn.

\begin{table}[h]
\centering
\caption{\textbf{Ablation: child-scorer conditioning.} Overall SYMSOL~I average
log-likelihood as conditioning components are removed.}
\label{tab:abl-cond}
\small
\begin{tabular}{l c}
\toprule
Conditioning & avg.\ log-lik. \\
\midrule
full (offset + parent + level) & \pend \\
\;\; $-$ child offset & \pend \\
\;\; $-$ parent rotation & \pend \\
\;\; $-$ level embedding & \pend \\
\;\; $-$ Fourier encoding (raw inputs) & \pend \\
\bottomrule
\end{tabular}
\end{table}

\paragraph{Training objective.}
A central claim is that the tree can be trained directly by exact maximum
likelihood, rather than as an energy with a sampled partition function. Keeping the
head fixed, \cref{tab:abl-obj} compares our tree-MLE objective against an InfoNCE
objective with importance-sampled negatives drawn from the tree (the SpyroPose-style
training), which SpyroPose reports trains more stably than a directly normalized
likelihood.

\begin{table}[h]
\centering
\caption{\textbf{Ablation: training objective.} Same head and backbone; exact
tree-MLE (ours) vs.\ an InfoNCE objective with tree-importance-sampled negatives.
Overall SYMSOL~I average log-likelihood, and negatives scored per step.}
\label{tab:abl-obj}
\small
\begin{tabular}{l c c}
\toprule
Objective & \#neg./step & avg.\ log-lik. \\
\midrule
InfoNCE + importance sampling & \pend & \pend \\
exact tree-MLE (ours) & 0 & \pend \\
\bottomrule
\end{tabular}
\end{table}

\section{\texorpdfstring{$\SEthree$}{SE(3)} construction and self-tests}
\label{sec:supp-se3}

The $\SEthree$ grid of \cref{sec:se3} is implemented as a product of the
rotation tree and a translation octree. The octree lives in a normalized unit cube
$g\in[-\tfrac12,\tfrac12]^3$, subdivided into the regular
$2^{l}\!\times\!2^{l}\!\times\!2^{l}$ grid; a cell is addressed by its integer
coordinates and, equivalently, by an octant path (octant $=b_x{+}2b_y{+}4b_z$,
most-significant bit first). The affine map $t=\hat{t}+Ag$ carries the cube into the
camera frame; since it is affine it preserves equal volume, so $V_{\text{box}}=|\!\det A|$
and each cell has volume $|\!\det A|/8^{l}$. A joint cell is the pair of a rotation
path and a translation path, and its per-node categorical is the $64$-way product
(child $=c_{\text{rot}}\!\cdot\!8+c_{\text{oct}}$) while both halves subdivide,
$8$-way afterwards, giving independent depths $L_r,L_t$.

Before any data or training we check four identities on the implementation.
(i)~Octree cell$\to$centre$\to$cell is exact at every level.
(ii)~A cell's octant path recomposes the same cell.
(iii)~A uniform predictor yields average log-likelihood exactly
$-\ln(\Vso\,V_{\text{box}})$ at every pair $(L_r,L_t)$, matching
\cref{eq:se3density} (maximum deviation $\sim\!10^{-15}$ nats over a range of
depths and a non-trivial affine box).
(iv)~The tree-structured softmax of empirical per-node conditionals equals the flat
empirical joint categorical to machine precision, confirming the factorization and
the density constant end to end without a network.
The box bounds come from a detection, as in~\cite{haugaard2023spyropose}; a pose
whose translation falls outside the box has infinite NLL, so we report the in-box
coverage rate alongside the likelihood.
\TODO{Report T-LESS/YCB-V numbers (\cref{tab:tless,tab:ycbv}) and the training
details for the $\SEthree$ head once the real-image runs complete.}
